% easybook — a pure academic template built on ctexbook, for Chinese
% theses or books. Ships with TeX Live; not bundled.
%
% Class documentation: texdoc easybook   |   https://ctan.org/pkg/easybook
% Licence: LPPL-1.3c, (c) 2021-2025 Qu Yi.
\documentclass[scheme=chinese,font=auto]{easybook}

\title{数据结构导论}
\author{王芳}
\date{\today}

\begin{document}

\frontmatter

\maketitle

\chapter*{前言}
本书是数据结构课程的配套讲义，覆盖线性表、栈、队列与二叉树的基本操作及其复杂度分析。每章配有示例代码与练习题。

\tableofcontents

\mainmatter

\chapter{线性表}

线性表是最基本的数据结构，元素之间存在线性的前后关系。常见实现方式有顺序表和链表两种。

\section{顺序表}

顺序表使用连续的内存空间存储元素，支持 $O(1)$ 的随机访问，但插入和删除操作在最坏情况下需要 $O(n)$ 时间，因为要移动后续元素。

\section{链表}

链表通过指针将不连续的节点串联起来。插入和删除操作只需 $O(1)$ 时间（在已知位置的情况下），但访问第 $k$ 个元素需要 $O(k)$ 时间。

\chapter{栈与队列}

栈遵循后进先出（LIFO）原则，队列遵循先进先出（FIFO）原则。两者都可以用顺序表或链表实现。

\section{栈的应用}

栈常用于表达式求值、括号匹配和函数调用的管理（调用栈）。

\section{队列的应用}

队列常用于广度优先搜索、任务调度和缓冲区管理。

\chapter{二叉树}

二叉树是每个节点最多有两个子节点的树形结构。二叉搜索树（BST）在平均情况下支持 $O(\log n)$ 的查找、插入和删除。

\end{document}
